\section{Related Work}

This work builds on the Cube Rule as a prior food-classification framework and introduces two external strands of literature to formalize the proposed tensor-based approach.
The first strand is modern LLM control and prompting: \term{prompt taxonomies}, \term{schema-constrained generation}, \term{grammar-constrained decoding}, \term{guided decoding}, and \term{retrieval-augmented generation} provide language for the prompt design, structured decoding, exemplar use, and post-generation filtering described in the Methods section.
The second strand is tensor and geometry-oriented representation theory: \term{tensor completion}, \term{manifold hypothesis}, \term{simplicial complex} constructions, \term{persistent homology}, and \term{topological data analysis} provide language for discussing empty-cell completion, clustered occupancy structure, and topological features of the support set $\mathcal{S} = \{(i,j,k) : T_{i,j,k}=1\}$ in higher-order categorical space.
Classical combinatorics references are also relevant for occupancy guarantees, including direct pigeonhole-style collision arguments when the number of foods exceeds $|I \times J \times K| = IJK$.
Taken together, these literatures provide the methodological and conceptual vocabulary used to define, populate, and analyze the proposed taxonomy.
